\IfFileExists{stacks-project.cls}{%
\documentclass{stacks-project}
}{%
\documentclass{amsart}
}

% For dealing with references we use the comment environment
\usepackage{verbatim}
\newenvironment{reference}{\comment}{\endcomment}
%\newenvironment{reference}{}{}
\newenvironment{slogan}{\comment}{\endcomment}
\newenvironment{history}{\comment}{\endcomment}

% For commutative diagrams we use Xy-pic
\usepackage[all]{xy}

% We use 2cell for 2-commutative diagrams.
\xyoption{2cell}
\UseAllTwocells

% We use multicol for the list of chapters between chapters
\usepackage{multicol}

% This is generally recommended for better output
\usepackage{lmodern}
\usepackage[T1]{fontenc}

% For cross-file-references
\usepackage{xr-hyper}

% Package for hypertext links:
\usepackage{hyperref}

% For any local file, say "hello.tex" you want to link to please
% use \externaldocument[hello-]{hello}
\externaldocument[introduction-]{introduction}
\externaldocument[conventions-]{conventions}
\externaldocument[sets-]{sets}
\externaldocument[categories-]{categories}
\externaldocument[topology-]{topology}
\externaldocument[sheaves-]{sheaves}
\externaldocument[sites-]{sites}
\externaldocument[stacks-]{stacks}
\externaldocument[fields-]{fields}
\externaldocument[algebra-]{algebra}
\externaldocument[brauer-]{brauer}
\externaldocument[homology-]{homology}
\externaldocument[derived-]{derived}
\externaldocument[simplicial-]{simplicial}
\externaldocument[more-algebra-]{more-algebra}
\externaldocument[smoothing-]{smoothing}
\externaldocument[modules-]{modules}
\externaldocument[sites-modules-]{sites-modules}
\externaldocument[injectives-]{injectives}
\externaldocument[cohomology-]{cohomology}
\externaldocument[sites-cohomology-]{sites-cohomology}
\externaldocument[dga-]{dga}
\externaldocument[dpa-]{dpa}
\externaldocument[sdga-]{sdga}
\externaldocument[hypercovering-]{hypercovering}
\externaldocument[schemes-]{schemes}
\externaldocument[constructions-]{constructions}
\externaldocument[properties-]{properties}
\externaldocument[morphisms-]{morphisms}
\externaldocument[coherent-]{coherent}
\externaldocument[divisors-]{divisors}
\externaldocument[limits-]{limits}
\externaldocument[varieties-]{varieties}
\externaldocument[topologies-]{topologies}
\externaldocument[descent-]{descent}
\externaldocument[perfect-]{perfect}
\externaldocument[more-morphisms-]{more-morphisms}
\externaldocument[flat-]{flat}
\externaldocument[groupoids-]{groupoids}
\externaldocument[more-groupoids-]{more-groupoids}
\externaldocument[etale-]{etale}
\externaldocument[chow-]{chow}
\externaldocument[intersection-]{intersection}
\externaldocument[pic-]{pic}
\externaldocument[weil-]{weil}
\externaldocument[adequate-]{adequate}
\externaldocument[dualizing-]{dualizing}
\externaldocument[duality-]{duality}
\externaldocument[discriminant-]{discriminant}
\externaldocument[derham-]{derham}
\externaldocument[local-cohomology-]{local-cohomology}
\externaldocument[algebraization-]{algebraization}
\externaldocument[curves-]{curves}
\externaldocument[resolve-]{resolve}
\externaldocument[models-]{models}
\externaldocument[functors-]{functors}
\externaldocument[equiv-]{equiv}
\externaldocument[pione-]{pione}
\externaldocument[etale-cohomology-]{etale-cohomology}
\externaldocument[proetale-]{proetale}
\externaldocument[relative-cycles-]{relative-cycles}
\externaldocument[more-etale-]{more-etale}
\externaldocument[trace-]{trace}
\externaldocument[crystalline-]{crystalline}
\externaldocument[spaces-]{spaces}
\externaldocument[spaces-properties-]{spaces-properties}
\externaldocument[spaces-morphisms-]{spaces-morphisms}
\externaldocument[decent-spaces-]{decent-spaces}
\externaldocument[spaces-cohomology-]{spaces-cohomology}
\externaldocument[spaces-limits-]{spaces-limits}
\externaldocument[spaces-divisors-]{spaces-divisors}
\externaldocument[spaces-over-fields-]{spaces-over-fields}
\externaldocument[spaces-topologies-]{spaces-topologies}
\externaldocument[spaces-descent-]{spaces-descent}
\externaldocument[spaces-perfect-]{spaces-perfect}
\externaldocument[spaces-more-morphisms-]{spaces-more-morphisms}
\externaldocument[spaces-flat-]{spaces-flat}
\externaldocument[spaces-groupoids-]{spaces-groupoids}
\externaldocument[spaces-more-groupoids-]{spaces-more-groupoids}
\externaldocument[bootstrap-]{bootstrap}
\externaldocument[spaces-pushouts-]{spaces-pushouts}
\externaldocument[spaces-chow-]{spaces-chow}
\externaldocument[groupoids-quotients-]{groupoids-quotients}
\externaldocument[spaces-more-cohomology-]{spaces-more-cohomology}
\externaldocument[spaces-simplicial-]{spaces-simplicial}
\externaldocument[spaces-duality-]{spaces-duality}
\externaldocument[formal-spaces-]{formal-spaces}
\externaldocument[restricted-]{restricted}
\externaldocument[spaces-resolve-]{spaces-resolve}
\externaldocument[formal-defos-]{formal-defos}
\externaldocument[defos-]{defos}
\externaldocument[cotangent-]{cotangent}
\externaldocument[examples-defos-]{examples-defos}
\externaldocument[algebraic-]{algebraic}
\externaldocument[examples-stacks-]{examples-stacks}
\externaldocument[stacks-sheaves-]{stacks-sheaves}
\externaldocument[criteria-]{criteria}
\externaldocument[artin-]{artin}
\externaldocument[quot-]{quot}
\externaldocument[stacks-properties-]{stacks-properties}
\externaldocument[stacks-morphisms-]{stacks-morphisms}
\externaldocument[stacks-limits-]{stacks-limits}
\externaldocument[stacks-cohomology-]{stacks-cohomology}
\externaldocument[stacks-perfect-]{stacks-perfect}
\externaldocument[stacks-introduction-]{stacks-introduction}
\externaldocument[stacks-more-morphisms-]{stacks-more-morphisms}
\externaldocument[stacks-geometry-]{stacks-geometry}
\externaldocument[moduli-]{moduli}
\externaldocument[moduli-curves-]{moduli-curves}
\externaldocument[examples-]{examples}
\externaldocument[exercises-]{exercises}
\externaldocument[guide-]{guide}
\externaldocument[desirables-]{desirables}
\externaldocument[coding-]{coding}
\externaldocument[obsolete-]{obsolete}
\externaldocument[fdl-]{fdl}
\externaldocument[index-]{index}

% Theorem environments.
%
\theoremstyle{plain}
\newtheorem{theorem}[subsection]{Theorem}
\newtheorem{proposition}[subsection]{Proposition}
\newtheorem{lemma}[subsection]{Lemma}

\theoremstyle{definition}
\newtheorem{definition}[subsection]{Definition}
\newtheorem{example}[subsection]{Example}
\newtheorem{exercise}[subsection]{Exercise}
\newtheorem{situation}[subsection]{Situation}

\theoremstyle{remark}
\newtheorem{remark}[subsection]{Remark}
\newtheorem{remarks}[subsection]{Remarks}

\numberwithin{equation}{subsection}

% Macros
%
\def\lim{\mathop{\mathrm{lim}}\nolimits}
\def\colim{\mathop{\mathrm{colim}}\nolimits}
\def\Spec{\mathop{\mathrm{Spec}}}
\def\Hom{\mathop{\mathrm{Hom}}\nolimits}
\def\Ext{\mathop{\mathrm{Ext}}\nolimits}
\def\SheafHom{\mathop{\mathcal{H}\!\mathit{om}}\nolimits}
\def\SheafExt{\mathop{\mathcal{E}\!\mathit{xt}}\nolimits}
\def\Sch{\mathit{Sch}}
\def\Mor{\mathop{\mathrm{Mor}}\nolimits}
\def\Ob{\mathop{\mathrm{Ob}}\nolimits}
\def\Sh{\mathop{\mathit{Sh}}\nolimits}
\def\NL{\mathop{N\!L}\nolimits}
\def\CH{\mathop{\mathrm{CH}}\nolimits}
\def\proetale{{pro\text{-}\acute{e}tale}}
\def\etale{{\acute{e}tale}}
\def\QCoh{\mathit{QCoh}}
\def\Ker{\mathop{\mathrm{Ker}}}
\def\Im{\mathop{\mathrm{Im}}}
\def\Coker{\mathop{\mathrm{Coker}}}
\def\Coim{\mathop{\mathrm{Coim}}}

% Boxtimes
%
\DeclareMathSymbol{\boxtimes}{\mathbin}{AMSa}{"02}

%
% Macros for moduli stacks/spaces
%
\def\QCohstack{\mathcal{QC}\!\mathit{oh}}
\def\Cohstack{\mathcal{C}\!\mathit{oh}}
\def\Spacesstack{\mathcal{S}\!\mathit{paces}}
\def\Quotfunctor{\mathrm{Quot}}
\def\Hilbfunctor{\mathrm{Hilb}}
\def\Curvesstack{\mathcal{C}\!\mathit{urves}}
\def\Polarizedstack{\mathcal{P}\!\mathit{olarized}}
\def\Complexesstack{\mathcal{C}\!\mathit{omplexes}}
% \Pic is the operator that assigns to X its picard group, usage \Pic(X)
% \Picardstack_{X/B} denotes the Picard stack of X over B
% \Picardfunctor_{X/B} denotes the Picard functor of X over B
\def\Pic{\mathop{\mathrm{Pic}}\nolimits}
\def\Picardstack{\mathcal{P}\!\mathit{ic}}
\def\Picardfunctor{\mathrm{Pic}}
\def\Deformationcategory{\mathcal{D}\!\mathit{ef}}
\usepackage{cleveref}
\crefname{originalsection}{original section}{original sections}
\Crefname{originalsection}{Original section}{Original sections}
\makeatletter
\newlabel{descent-section-descent-modules}{{023F}{source}{}{sourcepointer-1}{}}
\newlabel{descent-section-descent-modules@cref}{{[originalsection][0][]023F}{[1][0][]source}{}{}{}}
\makeatother
\crefname{originalreference}{original reference}{original references}
\Crefname{originalreference}{Original reference}{Original references}
\makeatletter
\newlabel{lemma-describe-sheaves-simplicial-site-site}{{09WF}{source}{}{sourcepointer-2}{}}
\newlabel{lemma-describe-sheaves-simplicial-site-site@cref}{{[originalreference][0][]09WF}{[1][0][]source}{}{}{}}
\makeatother
\crefname{originalremark}{original remark}{original remarks}
\Crefname{originalremark}{Original remark}{Original remarks}
\makeatletter
\newlabel{simplicial-remark-relations}{{016C}{source}{}{sourcepointer-3}{}}
\newlabel{simplicial-remark-relations@cref}{{[originalremark][0][]016C}{[1][0][]source}{}{}{}}
\makeatother
\crefname{originalreference}{original reference}{original references}
\Crefname{originalreference}{Original reference}{Original references}
\makeatletter
\newlabel{situation-simplicial-site}{{09WE}{source}{}{sourcepointer-4}{}}
\newlabel{situation-simplicial-site@cref}{{[originalreference][0][]09WE}{[1][0][]source}{}{}{}}
\makeatother
\title{Cartesian sheaves as descent data on a simplicial site}
\author{The Stacks Project Authors}
\date{stacks-project:spaces-simplicial.tex:local-snapshot-20261012; GNU FDL 1.2 or later}
\begin{document}
\maketitle
\noindent\textit{Editorial note.} All delivered mathematical text is the author's own, copied from the pinned Stacks Project chapter file sources/stacks-spaces-simplicial/pinned-original-source.tex; no statement or proof step was rewritten, reordered or removed.\par\medskip
\noindent\textit{Editorial note.} The counted proof is the author's complete argument: in-source Situation situation-simplicial-site; external stubs simplicial-remark-relations (Simplicial chapter), descent-section-descent-modules (Descent chapter).\par\medskip
\noindent\textit{Editorial note.} The label lemma-characterize-cartesian names the counted statement, and the proof environment immediately following it in the source is delivered verbatim as its proof.\par\medskip
\noindent\textit{Editorial note.} The counted statement is carried by the author's own lemma environment in the chapter source.\par\medskip
\noindent\textit{Allowed standard prerequisites (editorial).}\begin{itemize}
\item Simplicial face-map relations. Simplicial, Remark simplicial-remark-relations. the identities d\textasciicircum{}1\_1 d\textasciicircum{}2\_2 = d\textasciicircum{}1\_1 d\textasciicircum{}2\_1, d\textasciicircum{}1\_1 d\textasciicircum{}2\_0 = d\textasciicircum{}1\_0 d\textasciicircum{}2\_2, d\textasciicircum{}1\_0 d\textasciicircum{}2\_0 = d\textasciicircum{}1\_0 d\textasciicircum{}2\_1 that make the cocycle condition on alpha well defined.
\item Descent datum construction pattern. Descent, Section descent-section-descent-modules. the tau\textasciicircum{}n\_i and tau\textasciicircum{}n\_\{ij\} degeneracy maps and the quasi-inverse construction pattern borrowed for sheaves on the simplicial site.
\end{itemize}\medskip
\ifdefined\hypertarget\hypertarget{sourcepointer-1}{}\fi
\noindent\textit{Source reference.} Original section 023F: sources/stacks-chow/tags.txt, source lines 2740--2740. Prerequisite pointer invoked at proof line(s) 2613; it names a standard fact of a sibling chapter that the counted argument uses, and every step of the counted proof built on it is retained verbatim in the delivered body.\par\medskip
\ifdefined\hypertarget\hypertarget{sourcepointer-2}{}\fi
\noindent\textit{Source reference.} Original reference 09WF: sources/stacks-spaces-simplicial/pinned-original-source.tex, source lines 953--953. Prerequisite pointer invoked at proof line(s) 2637; it names a standard fact of the same chapter that the counted argument uses, and every step of the counted proof built on it is retained verbatim in the delivered body.\par\medskip
\ifdefined\hypertarget\hypertarget{sourcepointer-3}{}\fi
\noindent\textit{Source reference.} Original remark 016C: sources/stacks-chow/tags.txt, source lines 1539--1539. Prerequisite pointer invoked at proof line(s) 2589; it names a standard fact of a sibling chapter that the counted argument uses, and every step of the counted proof built on it is retained verbatim in the delivered body.\par\medskip
\ifdefined\hypertarget\hypertarget{sourcepointer-4}{}\fi
\noindent\textit{Source reference.} Original reference 09WE: sources/stacks-spaces-simplicial/pinned-original-source.tex, source lines 888--888. Prerequisite pointer invoked at proof line(s) 2562; it names a standard fact of the same chapter that the counted argument uses, and every step of the counted proof built on it is retained verbatim in the delivered body.\par\medskip
\begin{lemma}
\label{lemma-characterize-cartesian}
In Situation \ref{situation-simplicial-site}.
The category of cartesian sheaves of sets (resp.\ abelian groups)
is equivalent to the category of pairs $(\mathcal{F}, \alpha)$
where $\mathcal{F}$ is a sheaf of sets (resp.\ abelian groups)
on $\mathcal{C}_0$ and
$$
\alpha :
(f_{\delta_1^1})^{-1}\mathcal{F}
\longrightarrow (f_{\delta_0^1})^{-1}\mathcal{F}
$$
is an isomorphism of sheaves of sets (resp.\ abelian groups)
on $\mathcal{C}_1$ such that
$(f_{\delta^2_1})^{-1}\alpha =
(f_{\delta^2_0})^{-1}\alpha \circ (f_{\delta^2_2})^{-1}\alpha$
as maps of sheaves on $\mathcal{C}_2$.
\end{lemma}

\begin{proof}
We abbreviate
$d^n_j = f_{\delta^n_j} : \Sh(\mathcal{C}_n) \to \Sh(\mathcal{C}_{n - 1})$.
The condition on $\alpha$ in the statement of the lemma makes sense because
$$
d^1_1 \circ d^2_2 = d^1_1 \circ d^2_1, \quad
d^1_1 \circ d^2_0 = d^1_0 \circ d^2_2, \quad
d^1_0 \circ d^2_0 = d^1_0 \circ d^2_1
$$
as morphisms of topoi $\Sh(\mathcal{C}_2) \to \Sh(\mathcal{C}_0)$, see
Simplicial, Remark \ref{simplicial-remark-relations}. Hence we
can picture these maps as follows
$$
\xymatrix{
& (d^2_0)^{-1}(d^1_1)^{-1}\mathcal{F} \ar[r]_-{(d^2_0)^{-1}\alpha} &
(d^2_0)^{-1}(d^1_0)^{-1}\mathcal{F} \ar@{=}[rd] & \\
(d^2_2)^{-1}(d^1_0)^{-1}\mathcal{F} \ar@{=}[ru] & & &
(d^2_1)^{-1}(d^1_0)^{-1}\mathcal{F} \\
& (d^2_2)^{-1}(d^1_1)^{-1}\mathcal{F} \ar[lu]^{(d^2_2)^{-1}\alpha} \ar@{=}[r] &
(d^2_1)^{-1}(d^1_1)^{-1}\mathcal{F} \ar[ru]_{(d^2_1)^{-1}\alpha}
}
$$
and the condition signifies the diagram is commutative. It is clear that
given a cartesian sheaf $\mathcal{G}$ of sets (resp.\ abelian groups)
on $\mathcal{C}_{total}$
we can set $\mathcal{F} = \mathcal{G}_0$ and $\alpha$ equal to the composition
$$
(d_1^1)^{-1}\mathcal{G}_0 \to \mathcal{G}_1
\leftarrow (d_1^0)^{-1}\mathcal{G}_0
$$
where the arrows are invertible as $\mathcal{G}$ is cartesian.
To prove this functor
is an equivalence we construct a quasi-inverse. The construction of
the quasi-inverse is analogous to the construction discussed in
Descent, Section \ref{descent-section-descent-modules} from which we borrow
the notation $\tau^n_i : [0] \to [n]$, $0 \mapsto i$ and
$\tau^n_{ij} : [1] \to [n]$, $0 \mapsto i$, $1 \mapsto j$.
Namely, given a pair $(\mathcal{F}, \alpha)$
as in the lemma we set $\mathcal{G}_n = (f_{\tau^n_n})^{-1}\mathcal{F}$.
Given $\varphi : [n] \to [m]$ we define
$\mathcal{G}(\varphi) : (f_\varphi)^{-1}\mathcal{G}_n \to \mathcal{G}_m$
using
$$
\xymatrix{
(f_\varphi)^{-1}\mathcal{G}_n \ar@{=}[r] &
(f_\varphi)^{-1}(f_{\tau^n_n})^{-1}\mathcal{F} \ar@{=}[r] &
(f_{\tau^m_{\varphi(n)}})^{-1}\mathcal{F} \ar@{=}[r] &
(f_{\tau^m_{\varphi(n)m}})^{-1}(d^1_1)^{-1}\mathcal{F}
\ar[d]^{(f_{\tau^m_{\varphi(n)m}})^{-1}\alpha} \\
&
\mathcal{G}_m \ar@{=}[r] &
(f_{\tau^m_m})^{-1}\mathcal{F} \ar@{=}[r] &
(f_{\tau^m_{\varphi(n)m}})^{-1}(d^1_0)^{-1}\mathcal{F}
}
$$
We omit the verification that the commutativity of the displayed diagram
above implies the maps compose correctly and hence give rise to a
sheaf on $\mathcal{C}_{total}$, see
Lemma \ref{lemma-describe-sheaves-simplicial-site-site}.
We also omit the verification
that the two functors are quasi-inverse to each other.
\end{proof}

\end{document}
